\styledchapter{Aggregation}

How should a benevolent social planner decide the outcome if people report their individual preferences to her?

\section{Ordinal Aggregation}

\begin{notation}
	Let $R$ denote the set of weak preferences on $A$; i.e., the set of orderings, allowing ties.
\end{notation}

$R$ is the set of complete, transitive, and reflexive binary relations $\succsim$ on $A.$

\begin{definition}[Strict, weak social welfare function]
	A strict social welfare function is a function $F: P^N \to P$. A weak social welfare function is a function $F: P^N \to R$.
\end{definition}

Note that even weak social welfare functions are defined on the set of \emph{strict} preference profiles.

\begin{definition}[Kendall--Tau distance] \label{3-30-2}
	For any two strict rankings $\succ_i, \succ_i'$ on $A$, define the distance $d(\succ_i, \succ_i')$ between the two rankings as the number of pairs $a,b \in A$ such that $a \prec_i b$ and $a\succ_i' b$ or $a \succ_i b$ and $a\prec_i' b$.
\end{definition}

Note that the outcome of social welfare functions is a preference itself. Examples include:
\begin{enumerate}[label = (\roman*)]
	\item Dictatorship: $F\left( \succ_1, \cdots, \succ_{\left| N \right|} \right) = \succ_1$ for every $\left( \succ_1, \cdots, \succ_{\left| N \right|} \right)$.
	\item Elimination: rank in increasing order of round left.
	\item Borda count.
	\item Copeland method. 
	\item Kemeny--Young: Choose $F(\succ)$ to minimize \[
		\sum\limits_{i=1}^{\left| N \right|} d\left( F(\succ), \succ_i \right)
	.\] 
	\item Dodgson: For any outcome $a$ and preference profile $\succ$, let $\succ^{a}$ be the profile such that
		\begin{enumerate}[label = (\roman*)]
			\item $a$ is a Condorcet winner,
			\item $\sum_{i=1}^{\left| N \right|} d(\succ_i, \succ_i^a)$ is minimized.
		\end{enumerate}
	Then define $z(a) = \sum_{i=1}^{\left| N \right|} d(\succ_i, \succ_i^{a})$. Rank $a$ above $b$ iff $z(a) \le z(b)$.
\end{enumerate}

The elimination SWF is weak, as in a given round, the same number of people could rank the bottom two outcomes as their top choices. The Borda count is also weak.

\begin{definition}[Paretian property]
	A social welfare function $F$ satisfies the Paretian property if for any $\succ $ and $a,b \in A$ is such that if for all $i$, $a \succ_i b$, then \[
		a F(\succ) b
	.\] 
\end{definition}
\begin{definition}[Independence of irrelevant alternatives]
	A social welfare function $F$ satisfies independence of irrelevant alternatives (IIA) if for any two profiles $\succ,\succ'$ such that for all $i$, $\succ_i$ and $\succ_i'$ agree when restricted to $\left\{a,b\right\}$ for $a,b \in A$, then \[
		a F(\succ) b \iff a F(\succ') b
	.\] 
\end{definition}
	
\begin{example}[IIA assumption]
	Consider $A = \left\{a,b,c\right\}$ where $a \succ_1 b \succ_1 c$ and $b \succ_2 c \succ_2 a$. Without IIA, we might want to rank $b \succ_{S} a \succ_S c$, which seems sensible. Why would we want IIA? It is because in the preference $\succ_S$, we inferred cardinal rankings from ordinal rankings. To see why this is problematic, consider adding outcomes $d_1,\ldots,d_k$, where $d_\ell$ means that outcome $a$ is selected and you lose $\ell$ cents. Then
	\begin{align*}
		&a \succ_1' d_1 \succ_1' \cdots \succ_1' b \succ_1' c \\
		& b \succ_2' c \succ_2' a \succ_2' d_1 \succ_2' \cdots \succ_2' d_k
	,\end{align*}
	in which case IIA prohibits us from picking $a$.
\end{example}

The Borda count infers cardinal information when it shouldn't (and does not satisfy IIA); this is why in the Kendall--Tau distance, we just use the number of outcomes that the two preferences disagree on, and not anything relating to the ordering.

\begin{theorem}[Arrow's Impossibility]
	If $\left| A \right| > 2$, then the only (weak or strict) social welfare functions that satisfy IIA and the Paretian property are dictatorial.
\end{theorem}

\section{Cardinal Aggregation}

Assume that each agent $i$ has a utility function $u_i : A \to \R^+$, where we have a class $\mathscr{U}$ so $u_i \in \mathscr{U}$. We want to provide axioms on a function $F$ mapping $\mathscr{U}^{N} \to R$ or $\mathscr{U}^{N} \to U$.
For each $a \in A$, we can stack the utilities into a vector $\left( u_1(a), u_2(a), \ldots,u_{\left| N \right|}(a) \right)'$. Instead of defining axioms on $F$, we could provide axioms on a binary relation $\succsim_s$ on $\R_+^{\left| N \right|}$. This would allow us to compare outcomes $a,b \in A$ by comparing their vectors $\left( u_1(a),\ldots,u_{\left| N \right|}(a) \right)$ and $\left( u_1(b),\ldots,u_{\left| N \right|}(b) \right)$.
We could use this to form a function $F: \mathscr{U}^{\left| N \right|} \to R$, where \[
	a \succsim b \iff \left( u_1(a),\ldots,u_{\left| N \right|}(a) \right) \succsim \left( u_1(b),\ldots,u_{\left| N \right|}(b) \right)
.\] 

\begin{definition}[Cardinal social welfare ranking]
	A cardinal social welfare ranking is a complete, transitive and reflexive binary relation $\succsim_s$ on $\R_+^{\left| N \right|}$.
\end{definition}

\begin{example}[Cardinal social welfare rankings]
	We have some examples:\footnote{The leximin ranking is discussed further later in this section.}
	\begin{enumerate}[label = (\roman*)]
		\item Utilitarian: $u \succsim_s u'$ iff $\sum_i u_i \ge \sum_i u_i'$ 
		\item Rawlsian: $u \succsim_s u'$ iff $\min_i u_i \ge \min_i u_i'$ 
		\item Leximin: $u \succ_s v \iff u_k^* \succ v_k^*,$  where $u^*,v^*$ are $u,v$ rearranged in increasing order, and $k$ is the first index where $u,v$ differ.
	\end{enumerate}
\end{example}

Now we define properties we want cardinal rankings to satisfy.
\begin{definition}[Continuous cardinal social welfare ranking]
	We say that $\succsim_s$ is continuous if for any $u \in \R_+^{\left| N \right|}$, the sets \[
		\left\{u': u' \succsim_s u\right\}, \quad \left\{u' : u \succsim_s u'\right\}
	\]  are closed. 
	In other words, if $\left\{u^n\right\} \to u^*$ is such that $u^n \succsim_s u'$ for each $n\in \N$, then $u^* \succsim_s u'$. Similarly, if $u' \succsim_s u^n$ for each $n\in \N$, then $u' \succsim_s u^*$.
\end{definition}

\begin{definition}[Independence of unconcerned agents]
	$\succsim_s$ is independent of unconcerned agents if for any $x,y \in \R_+$ and any agent $i$, we have \[
		(x, u_{-i}) \succsim_s (x, u_{-i}') \iff (y,u_{-i}) \succsim_s (y,u_{-i}')
	\] 
	for any $u_{-i}, u_{-i}'$.
\end{definition}

We can represent a well-behaved ranking by a function that is additively separable across agents.

\begin{theorem}[Debreu]
	For any $\succsim_s$ satisfying continuity and independence of unconcerned agents, there exist functions $\left\{g_i\right\}_{i \in N}$ such that \[
		u\succsim_s u' \iff \sum\limits_{i}^{} g_i(u_i) \ge \sum\limits_{i}^{} g_i(u_i')
	.\] 
\end{theorem}

We have the following analogue of the Paretian property.

\begin{definition}[Strict monotonicity]
	We say that $\succsim_s$ is strictly monotonic if whenever $u_i \ge u_i'$ for all $i$ and there is some $j$ with $u_j > u_j'$, then $u \succ_s u'$.
\end{definition}

\begin{definition}[Scale invariance]
	We say $\succsim_s$ is scale invariant if \[
		u \succsim_s u' \iff \lambda u \succsim_s \lambda u', \quad \forall\  \lambda \in \R_{>0}
	.\] 
\end{definition}

\begin{theorem} \label{4-1-1}
	If $\succsim_s$ satisfies continuity, independence of unconcerned agents, strict monotonicity, and scale invariance, and if $\left| N \right| \ge 2$, then there are positive weights $w_{1},\ldots,w_{\left| N \right|}$ which sum to $1$ and some $\gamma \in \R$ such that \[
		u \succsim_s v \iff \sum\limits_{i}^{} w_i \frac{u_i^{1-\gamma}-1}{1-\gamma} \ge \sum\limits_{i}^{} w_i \frac{v_i^{1-\gamma} - 1}{1- \gamma}
	.\] 
\end{theorem}

We now explore properties of this CRRA transformation of the utilities. We see that when $\gamma=0$, we have \[
	u \succsim_s v \iff \sum\limits_{i}^{} w_i u_i \ge \sum\limits_{i}^{} w_i v_i
,\] which is a weighted utilitarian preference.

Furthermore, \[
	\lim\limits_{\gamma \uparrow 1} \frac{u^{1-\gamma}-1}{1-\gamma} = \ln{\left(u\right)}
,\] so for $\gamma$ close to $1$,  \[
	u\succsim_s v  \iff \sum\limits_{i}^{} w_i \ln{\left(u_i\right)} \ge \sum\limits_{i}^{} w_i \ln{\left(v_i\right)} \iff \prod\limits_{i}^{} u_i^{w_i} \ge \prod\limits_{i}^{}  v_i ^{w_i}
,\] 
which is called \emph{Nash welfare}.

As $\gamma\to \infty$, we get the leximin\footnote{A portmanteau of lexicographic and min.} order, in that \[
	u \succ_s v \iff u_k^* \succ v_k^*
,\] 
where $u^*,v^*$ are $u,v$ rearranged in increasing order, and $k$ is the first index where $u,v$ differ. The leximin ordering is a refinement (in the English, not mathematical sense) of the Rawlsian ordering.

Note that Theorem (\ref{4-1-1}) does not say that we can pick $\gamma$; instead, after fixing some ranking $\succsim_s$ of the planner, then we know there is \emph{some} $\gamma$ that represents the planner's rankings in this additively separable, CRRA way.

\section{Aggregation Under Uncertainty}

Suppose there is a finite number of states $\Theta = \left\{\theta_1,\ldots,\theta_m\right\}$. Each agent $i$ has utility $u_i: \Theta \to \R$ and evaluates a lottery $p$ over $\Theta$ by the Von Neumann--Morgenstern expected utility function \[
	U_i(p) = \sum\limits_{k=1}^{m} p_k \underbrace{u_i (\theta_k)}_{\mathclap{\text{Bernoulli index}}}
.\] 
(We use $U$ to denote utility over $\Delta(\Theta)$, including degenerate lotteries, and $u$ to denote utility over $\Theta$.)

There are good normative reasons for people to be VNM maximizers, but in an experimental setting, people violate this utility structure.
We assume that the planner is also VNM rational with utility \[
	U_s(p) = \sum\limits_{k=1}^{m} p_k u_s(\theta_k)
.\] 

\begin{definition}[Paretian property]
	The planner's utility function $U_s(\cdot)$ satisfies the Paretian property if $U_i(p) \ge U_i(q)$ for all $i$ implies $U_s(p) \ge U_s(q)$.
\end{definition}

\begin{theorem}[Harsanyi]
	The planner's utility $U_s$ is Paretian iff there is a constant $\beta$ and weights $\left( w_i \right)_{i \in N}$ which are nonnegative such that \[
		u_s(\theta) = \beta + \sum\limits_{i}^{} w_i u_i(\theta), \quad \forall\  \theta \in \Theta
	.\] Equivalently, we have a necessary and sufficient condition for the Paretian property of \[
		U_s(p) = \beta + \sum\limits_{i}^{} w_i U_i(p)
	.\] 
\end{theorem}

To prove this theorem, we first recall Lemma (\ref{4-1-2}). 

\begin{proof}
	We can write this as a linear programming problem, where the theorem holds if there exist $\beta$ and $(w_i)_{i \in N}$ such that \[
		\begin{pmatrix}
			1 & u_1(\theta_1) & u_2(\theta_1) & \cdots & u_{\left| N \right|}(\theta_1) \\
			1 & u_1(\theta_2) & u_2(\theta_2) & \cdots & u_{\left| N \right|}(\theta_2) \\
			\vdots & \vdots & \vdots & \ddots & \vdots \\
			1 & u_1(\theta_{m}) & u_2(\theta_m) & \cdots & u_{\left| N \right|}(\theta_m)
		\end{pmatrix} \begin{pmatrix}
		\beta \\
		w_1 \\
		w_2 \\
		\vdots \\
		w_{\left| N \right|}
		\end{pmatrix} = \begin{pmatrix}
		u_s(\theta_1) \\
		u_s(\theta_2) \\
		\vdots \\
		u_{s}(\theta_m)
		\end{pmatrix} 
	\] and \[
	-\begin{pmatrix}
		0 & 1 & 0 & \cdots & 0 \\
		0 & 0 & 1 & \cdots & 0 \\
		0 &  &  & \ddots & 0 \\
		0 & 0 & 0 & \cdots & 1
	\end{pmatrix} \begin{pmatrix}
		\beta \\
		w_1 \\
		w_2 \\
		\vdots \\
		w_{\left| N \right|}
		\end{pmatrix} \le \begin{pmatrix}
		0  \\
		0 \\ \vdots \\ 0
		\end{pmatrix} 
	.\] 

	If the feasibility alternative in Farkas' Lemma holds, then we are done. So, assume that we do not have satisfactory $\beta$ and $(w_i)_{i \in N}$.
	To apply Lemma (\ref{4-1-2}), take the vector of unknowns to be
	\[
		x = \left(\beta,w_1,\ldots,w_{\left| N \right|}\right)'.
	\]
	The first displayed system above is the equality system $Bx=c$, with one equality for each state $\theta_k$. The second displayed system is the inequality system $Ax \le 0$, with one inequality $-w_j \le 0$ for each agent $j$.
	Hence the Farkas certificate has two kinds of multipliers. The equality constraints receive unrestricted multipliers $q_1,\ldots,q_m$, while the inequalities $-w_j \le 0$ receive nonnegative multipliers $p_1,\ldots,p_{\left| N \right|}$.
	Writing $p$ and $q$ as row vectors, the certificate alternative in Lemma (\ref{4-1-2}) says that, if no such $\beta,w$ exist, then these multipliers must satisfy
	\[
		pA + qB = 0, \qquad 0 > p \cdot b + q \cdot c = p \cdot 0 + q\cdot c = q \cdot c
	.\]
	Here, the inequality matrix $A$ has rows
	\[
		\left(0, -e_j\right), \qquad j \in N,
	\]
	where the $0$ is the coefficient on $\beta$ and $-e_j$ selects the coefficient on $w_j$. Hence
	\[
		pA = \left(0,-p_1,\ldots,-p_{\left| N \right|}\right).
	\]
	The equality matrix $B$ has one row for each state:
	\[
		\left(1,u_1(\theta_k),\ldots,u_{\left| N \right|}(\theta_k)\right), \qquad k=1,\ldots,m.
	\]
	Therefore
	\[
		qB =
		\left(
			\sum\limits_{k=1}^{m} q_k,
			\sum\limits_{k=1}^{m} q_k u_1(\theta_k),
			\ldots,
			\sum\limits_{k=1}^{m} q_k u_{\left| N \right|}(\theta_k)
		\right).
	\]
	The condition $pA+qB=0$ means that each component of this row vector is zero. From the first component, which is the coefficient on $\beta$, we have
	\begin{equation}
		\sum\limits_{k=1}^{m} q_k = 0 \label{eq:4-1-1}
	,\end{equation}
	while the component corresponding to $w_j$ is
	\[
		-p_j + \sum\limits_{k=1}^{m} q_k u_j(\theta_k),
	\]
	so the coefficients on the $w_j$ imply
	\begin{equation}
		\sum\limits_{k=1}^{m} q_k u_j(\theta_k)-p_j = 0, \quad \forall\ j \label{eq:4-1-2}
	.\end{equation}
	Lastly, $c$ is the right-hand side of the equality system, so
	\[
		c =
		\left(u_s(\theta_1),\ldots,u_s(\theta_m)\right)'.
	\]
	Thus the strict inequality $q \cdot c < 0$ is
	\begin{equation}
		\sum\limits_{k=1}^{m} q_k u_s(\theta_k) < 0 \label{eq:4-1-3}
	.\end{equation}

	Because $p_j \ge 0$, Equation (\ref{eq:4-1-2}) implies
	\[
		\sum\limits_{k=1}^{m} q_k u_j(\theta_k) \ge 0, \quad \forall\ j
	.\] 
	Notice that if $q$ solves the above three constraints, then for every $\lambda > 0$, $\lambda q$ is a solution as well. Rescale $q$ such that \[
		\left| q_k \right| \le \frac{1}{m}
	\] for each $k$ and define $r_k = q_k + \frac{1}{m}$ for all $k$. 

	Returning to the three constraints with $r$ replacing $q$, 
	\begin{align*}
		\sum\limits_{k=1}^{m} r_k &= 1, \quad r_k \ge 0, \qquad \forall\ k \\
		\sum\limits_{k=1}^{m} r_k u_j(\theta_k) - \sum\limits_{k=1}^{m} \frac{1}{m} u_j(\theta_k) &\ge 0, \quad \forall\ j \\
		\sum\limits_{k=1}^{m} r_k u_s(\theta_k) - \sum\limits_{k=1}^{m} \frac{1}{m} u_s(\theta_k) &< 0
	.\end{align*}
	The second line says that every agent $j$ weakly prefers the lottery defined by $(r_k)$ to the uniform lottery, while the planner strictly prefers the uniform lottery. This is a violation of the Paretian property.
\end{proof}

